% GENERATED FILE. Edit canonical lessons, then run npm run generate:books.
% canonical-source-hash: fnv1a64:9f7ee8c8c49b58d0
% canonical-lessons: TA-C336-vilttu, TA-C336-kaipparru, TA-C336-caranatai, TA-C336-atipani, TA-C336-akkirami

\chapter{Victory and Defeat}
\label{ch:ta-victory-and-defeat}
\hlchaptermodality{336}

\begin{chapteropening}
\textbf{By the end of this chapter:} \emph{I can say how a war was won and lost.}

Complete the last lesson of chapter 336: I can say how a war was won and lost.
\end{chapteropening}

\section[\texorpdfstring{\ta{வீழ்த்து} (vīḻttu)}{வீழ்த்து (vīḻttu)}]{\ta{வீழ்த்து} (vīḻttu) — to bring down; to topple}

\label{lesson:TA-C336-vilttu}



\begin{quote}
Before the new one: say the Tamil for a fortification, then the Tamil for a spy.
\end{quote}

\subsection*{You'll want to know: \ta{வீழ்த்து}}
\textbf{\ta{வீழ்த்து}} — \emph{vīḻttu} — \textquotedblleft{}to bring down; to topple\textquotedblright{}.

\begin{quote}
\ta{சிறிய} \ta{படை} \ta{பெரிய} \ta{படையை} \ta{வீழ்த்தியது}.
\end{quote}

\begin{quote}
\emph{ciṟiya paṭai periya paṭaiyai vīḻttiyatu.}
\end{quote}

\begin{quote}
The small army brought down the big army.
\end{quote}

Past (I …): \ta{வீழ்த்தினேன்}, \emph{vīḻttiṉēṉ}.

Causative of \ta{வீழ்}, to fall: to make fall.

\begin{cousinweb}
\textbf{\ta{வீழ்த்து}} (\emph{vīḻttu}) is native Dravidian: DEDR 5430 gives it as the causative of \ta{வீழ்}, \ta{விழு}, \textquotedblleft{}to fall\textquotedblright{}.
\end{cousinweb}

\begin{grammarlens}[title={What you've built}]
One more word for telling what happened.
\end{grammarlens}

\subsection*{Your turn}
\begin{itemize}
\raggedright
  \item \emph{Say it:} \emph{vīḻttu}
  \item \emph{Say it:} \emph{vīḻttu}, once more
  \item \emph{Say it:} \emph{ciṟiya paṭai periya paṭaiyai vīḻttiyatu.} — The small army brought down the big army.
  \item \emph{Recall:} say the Tamil for a fortification, then the Tamil for a spy, then say \emph{vīḻttu} again
\end{itemize}

\begin{tcolorbox}[breakable,colback=teal!4,colframe=teal!35!black,title={Before you move on}]
What does \ta{வீழ்த்து} mean? (\textquotedblleft{}To bring down; to topple\textquotedblright{}.) Say it once more.
\end{tcolorbox}
\section[\texorpdfstring{\ta{கைப்பற்று} (kaippaṟṟu)}{கைப்பற்று (kaippaṟṟu)}]{\ta{கைப்பற்று} (kaippaṟṟu) — to capture; to seize}

\label{lesson:TA-C336-kaipparru}



\begin{quote}
Before the new one: say the Tamil for a spy, then the Tamil for to bring down.
\end{quote}

\subsection*{You'll want to know: \ta{கைப்பற்று}}
\textbf{\ta{கைப்பற்று}} — \emph{kaippaṟṟu} — \textquotedblleft{}to capture; to seize\textquotedblright{}.

\begin{quote}
\ta{படை} \ta{மூன்று} \ta{நாளில்} \ta{கோட்டையைக்} \ta{கைப்பற்றியது}.
\end{quote}

\begin{quote}
\emph{paṭai mūṉṟu nāḷil kōṭṭaiyaik kaippaṟṟiyatu.}
\end{quote}

\begin{quote}
The army captured the fort in three days.
\end{quote}

Past (I …): \ta{கைப்பற்றினேன்}, \emph{kaippaṟṟiṉēṉ}.

Also of goods: \ta{காவலர்கள்} \ta{பணத்தைக்} \ta{கைப்பற்றினார்கள்}.

\begin{cousinweb}
\textbf{\ta{கைப்பற்று}} (\emph{kaippaṟṟu}) is a Tamil compound built from \ta{கை}, \textquotedblleft{}hand\textquotedblright{} (DEDR 2023), and \ta{பற்று}, \textquotedblleft{}to grasp, seize\textquotedblright{} (DEDR 4034).
\end{cousinweb}

\begin{grammarlens}[title={What you've built}]
One more word for telling what happened.
\end{grammarlens}

\subsection*{Your turn}
\begin{itemize}
\raggedright
  \item \emph{Say it:} \emph{kaippaṟṟu}
  \item \emph{Say it:} \emph{kaippaṟṟu}, once more
  \item \emph{Say it:} \emph{paṭai mūṉṟu nāḷil kōṭṭaiyaik kaippaṟṟiyatu.} — The army captured the fort in three days.
  \item \emph{Recall:} say the Tamil for a spy, then the Tamil for to bring down, then say \emph{kaippaṟṟu} again
\end{itemize}

\begin{tcolorbox}[breakable,colback=teal!4,colframe=teal!35!black,title={Before you move on}]
What does \ta{கைப்பற்று} mean? (\textquotedblleft{}To capture; to seize\textquotedblright{}.) Say it once more.
\end{tcolorbox}
\section[\texorpdfstring{\ta{சரணடை} (caraṇaṭai)}{சரணடை (caraṇaṭai)}]{\ta{சரணடை} (caraṇaṭai) — to surrender}

\label{lesson:TA-C336-caranatai}



\begin{quote}
Before the new one: say the Tamil for to bring down, then the Tamil for to capture.
\end{quote}

\subsection*{You'll want to know: \ta{சரணடை}}
\textbf{\ta{சரணடை}} — \emph{caraṇaṭai} — \textquotedblleft{}to surrender\textquotedblright{}.

\begin{quote}
\ta{இறுதியில்} \ta{கோட்டையின்} \ta{வீரர்கள்} \ta{சரணடைந்தார்கள்}.
\end{quote}

\begin{quote}
\emph{iṟutiyil kōṭṭaiyiṉ vīrarkaḷ caraṇaṭaintārkaḷ.}
\end{quote}

\begin{quote}
In the end the fort's soldiers surrendered.
\end{quote}

Past (I …): \ta{சரணடைந்தேன்}, \emph{caraṇaṭaintēṉ}.

To whom: -\ta{இடம்}, \ta{மன்னரிடம்} \ta{சரணடை}.

\begin{cousinweb}
\textbf{\ta{சரணடை}} (\emph{caraṇaṭai}) joins \ta{சரண்}, \textquotedblleft{}refuge\textquotedblright{}, borrowed from Sanskrit \emph{śaraṇa}, \textquotedblleft{}shelter\textquotedblright{} (CDIAL 12326), with native \ta{அடை}, \textquotedblleft{}to reach, take refuge in\textquotedblright{} (DEDR 79).
\end{cousinweb}

\begin{grammarlens}[title={What you've built}]
One more word for telling what happened.
\end{grammarlens}

\subsection*{Your turn}
\begin{itemize}
\raggedright
  \item \emph{Say it:} \emph{caraṇaṭai}
  \item \emph{Say it:} \emph{caraṇaṭai}, once more
  \item \emph{Say it:} \emph{iṟutiyil kōṭṭaiyiṉ vīrarkaḷ caraṇaṭaintārkaḷ.} — In the end the fort's soldiers surrendered.
  \item \emph{Recall:} say the Tamil for to bring down, then the Tamil for to capture, then say \emph{caraṇaṭai} again
\end{itemize}

\begin{tcolorbox}[breakable,colback=teal!4,colframe=teal!35!black,title={Before you move on}]
What does \ta{சரணடை} mean? (\textquotedblleft{}To surrender\textquotedblright{}.) Say it once more.
\end{tcolorbox}
\section[\texorpdfstring{\ta{அடிபணி} (aṭipaṇi)}{அடிபணி (aṭipaṇi)}]{\ta{அடிபணி} (aṭipaṇi) — to submit; to bow down}

\label{lesson:TA-C336-atipani}



\begin{quote}
Before the new one: say the Tamil for to capture, then the Tamil for to surrender.
\end{quote}

\subsection*{You'll want to know: \ta{அடிபணி}}
\textbf{\ta{அடிபணி}} — \emph{aṭipaṇi} — \textquotedblleft{}to submit; to bow down\textquotedblright{}.

\begin{quote}
\ta{சிறிய} \ta{நாடுகள்} \ta{பேரரசுக்கு} \ta{அடிபணிந்தன}.
\end{quote}

\begin{quote}
\emph{ciṟiya nāṭukaḷ pēraracukku aṭipaṇintaṉa.}
\end{quote}

\begin{quote}
The small kingdoms submitted to the empire.
\end{quote}

Past (I …): \ta{அடிபணிந்தேன்}, \emph{aṭipaṇintēṉ}.

To whom: -\ta{க்கு}; \ta{அடி}, foot + \ta{பணி}, to bow.

\begin{cousinweb}
\textbf{\ta{அடிபணி}} (\emph{aṭipaṇi}) is a Tamil compound built from \ta{அடி}, \textquotedblleft{}foot\textquotedblright{} (DEDR 72), and \ta{பணி}, \textquotedblleft{}to be lowered, be humble\textquotedblright{} (DEDR 3888).
\end{cousinweb}

\begin{grammarlens}[title={What you've built}]
One more word for telling what happened.
\end{grammarlens}

\subsection*{Your turn}
\begin{itemize}
\raggedright
  \item \emph{Say it:} \emph{aṭipaṇi}
  \item \emph{Say it:} \emph{aṭipaṇi}, once more
  \item \emph{Say it:} \emph{ciṟiya nāṭukaḷ pēraracukku aṭipaṇintaṉa.} — The small kingdoms submitted to the empire.
  \item \emph{Recall:} say the Tamil for to capture, then the Tamil for to surrender, then say \emph{aṭipaṇi} again
\end{itemize}

\begin{tcolorbox}[breakable,colback=teal!4,colframe=teal!35!black,title={Before you move on}]
What does \ta{அடிபணி} mean? (\textquotedblleft{}To submit; to bow down\textquotedblright{}.) Say it once more.
\end{tcolorbox}
\section[\texorpdfstring{\ta{ஆக்கிரமி} (ākkirami)}{ஆக்கிரமி (ākkirami)}]{\ta{ஆக்கிரமி} (ākkirami) — to occupy; to invade}

\label{lesson:TA-C336-akkirami}



\begin{quote}
Before the new one: say the Tamil for to surrender, then the Tamil for to submit.
\end{quote}

\subsection*{You'll want to know: \ta{ஆக்கிரமி}}
\textbf{\ta{ஆக்கிரமி}} — \emph{ākkirami} — \textquotedblleft{}to occupy; to invade\textquotedblright{}.

\begin{quote}
\ta{அந்தப்} \ta{படை} \ta{பக்கத்து} \ta{நாட்டை} \ta{ஆக்கிரமித்தது}.
\end{quote}

\begin{quote}
\emph{antap paṭai pakkattu nāṭṭai ākkiramittatu.}
\end{quote}

\begin{quote}
That army occupied the neighbouring country.
\end{quote}

Past (I …): \ta{ஆக்கிரமித்தேன்}, \emph{ākkiramittēṉ}.

Also of land taken unlawfully: \ta{நிலத்தை} \ta{ஆக்கிரமி}.

\begin{cousinweb}
\textbf{\ta{ஆக்கிரமி}} (\emph{ākkirami}) is a Tamil verb built on Sanskrit \emph{ākrama}, \textquotedblleft{}approaching, overcoming\textquotedblright{} (Monier-Williams).
\end{cousinweb}

\begin{grammarlens}[title={What you've built}]
One more word for telling what happened.
\end{grammarlens}

\subsection*{Your turn}
\begin{itemize}
\raggedright
  \item \emph{Say it:} \emph{ākkirami}
  \item \emph{Say it:} \emph{ākkirami}, once more
  \item \emph{Say it:} \emph{antap paṭai pakkattu nāṭṭai ākkiramittatu.} — That army occupied the neighbouring country.
  \item \emph{Recall:} say the Tamil for to surrender, then the Tamil for to submit, then say \emph{ākkirami} again
\end{itemize}

\begin{tcolorbox}[breakable,colback=teal!4,colframe=teal!35!black,title={Before you move on}]
What does \ta{ஆக்கிரமி} mean? (\textquotedblleft{}To occupy; to invade\textquotedblright{}.) Say it once more.
\end{tcolorbox}
